\section{Gauge-covariant response and statistical estimators}
\label{app:methods}

This section records the projector identities and statistical contractions
used in the main text.  They are stated independently of any particular parent
Hamiltonian so that model-dependent exactness and gap checks can be audited
separately from the response calculation.

\subsection{Differentiation of an isolated spectral projector}

Assume that $H(\boldsymbol\lambda)$ is $C^1$ in a neighborhood of the base
point and has a rank-$D$ eigenprojector $P$ separated from $Q=1-P$ by
$\Delta>0$.  Differentiating $P^2=P$ gives
\[
 P(\partial_aP)P=0,
 \qquad
 Q(\partial_aP)Q=0,
\]
so $\partial_aP$ is off diagonal in the $P/Q$ decomposition.  Differentiating
$HP=E_0P$, multiplying by $Q$ on the left and $P$ on the right, and using
$QHP=0$ yields
\[
 Q(H-E_0)Q\,Q(\partial_aP)P=-Q(\partial_aH)P.
\]
The inverse on $Q\cH$ gives
\begin{equation}
 X_a=Q(\partial_aP)P
 =-Q\bigl[Q(H-E_0)Q\bigr]^{-1}Q(\partial_aH)P.
 \label{eq:supp-response}
\end{equation}
Hermiticity supplies the
other block, $P(\partial_aP)Q=X_a^\dagger$, hence
\[
 \partial_aP=X_a+X_a^\dagger.
\]
The norm estimate
$\lVert X_a\rVert\leq\lVert Q(\partial_aH)P\rVert/\Delta$
is useful as a conditioning check, not as a prediction for the distribution
of response matrix elements.

The same identities follow from the Riesz formula
\[
 P=-\frac{1}{2\pi i}\oint_\Gamma(H-z)^{-1}\,dz,
\]
where $\Gamma$ encloses only $E_0$.  This representation makes clear that the
construction is independent of a basis chosen inside the degenerate fiber.
It also states the failure mode precisely: if no contour separates the target
cluster from the remaining spectrum, a fixed-rank adiabatic projector is no
longer defined by this prescription.

\subsection{Frame transformations and closed response words}

In a local frame $\Psi$, the rectangular response is
$x_a=Q(\partial_aP)\Psi$.  Under $\Psi\mapsto\Psi U$,
\[
 x_a\mapsto x_aU,
 \qquad
 x_a^\dagger x_b\mapsto U^\dagger x_a^\dagger x_bU.
\]
Consequently, a trace word is independent of the frame whenever its fiber
indices form a closed cycle.  Examples include
\[
 \Tr(x_a^\dagger x_b),\qquad
 \Tr(x_a^\dagger x_bx_c^\dagger x_d),\qquad
 \Tr(F_{ab}F_{cd}).
\]
An individual entry of $x_a$ is covariant and can be used inside an estimator
whose contractions transform consistently.  It is not, on its own, a
gauge-invariant measurement.  An individual entry of the connection is even
less suitable because of the inhomogeneous derivative term in its gauge
transformation.

For ensembles at a common base projector, one frame is fixed before sampling
the tangent panel; a subsequent frame change is applied to every panel.  For
independent Hamiltonians, there is generally no canonical identification of
their protected fibers.  Such ensembles are compared through invariant scalar
statistics unless a parallel-transport or polar-overlap identification is
specified before the data are opened.

\subsection{Complete complex Wick tensor}

Let $(X_a)_{\alpha i}$ denote the response from fiber index $i$ to complement
index $\alpha$, after centering and the registered channel whitening.  Define
the full entry covariance and pseudocovariance by
\[
 C_{ab;\alpha i,\beta j}
 =\mathbb E\bigl[(X_a)_{\alpha i}(X_b)_{\beta j}^{*}\bigr],
 \qquad
\Pi_{ab;\alpha i,\beta j}
 =\mathbb E\bigl[(X_a)_{\alpha i}(X_b)_{\beta j}\bigr].
\]
No tensor-product decomposition of these arrays is assumed.  The physical
four-point tensor in these indices is
\begin{equation}
 T_{abcd}=\frac{1}{D}\sum_{\alpha,\beta,i,j}
 \mathbb E\bigl[
 (X_a)_{\alpha i}^{*}(X_b)_{\alpha j}
 (X_c)_{\beta j}^{*}(X_d)_{\beta i}
 \bigr].
 \label{eq:supp-fourpoint}
\end{equation}
Applying the three Wick pairings to Eq.~\eqref{eq:supp-fourpoint} gives
\begin{align}
 W_{abcd}=\frac{1}{D}\sum_{\alpha,\beta,i,j}\Bigl[{}
 &C_{ba;\alpha j,\alpha i}C_{dc;\beta i,\beta j}
 \nonumber\\
 &+\Pi_{ac;\alpha i,\beta j}^{*}\Pi_{bd;\alpha j,\beta i}
 \nonumber\\
 &+C_{da;\beta i,\alpha i}C_{bc;\alpha j,\beta j}
 \Bigr].
 \label{eq:supp-wick}
\end{align}
The first and third lines pair one conjugated entry with one unconjugated
entry.  The middle line pairs the two conjugated entries with each other and
the two unconjugated entries with each other; it is exactly the contribution
controlled by the pseudocovariance.  Equation~\eqref{eq:supp-wick} is therefore
valid for a general zero-mean complex Gaussian vector.  The two-pairing formula
is recovered only after the additional condition $\Pi=0$ has been established.

The augmented covariance
\[
 \Gamma=
 \begin{pmatrix}
 C&\Pi\\
 \Pi^*&C^*
 \end{pmatrix}
\]
must be positive semidefinite.  Complete whitening can be implemented as a
real-linear transformation on the support of $\Gamma$.  Equivalently, one can
retain the original response coordinates and subtract
$W(C,\Pi)$.  We use the latter representation because it keeps the physical
tangent labels visible.  Exact null directions are removed according to a
rank tolerance fixed by the second-moment audit, never by minimizing the
fourth-order residual.

\subsection{Pair-split estimator and jackknife uncertainty}

Let $r=1,\ldots,R$ label the declared statistical units.  For a
fixed-Hamiltonian tangent-panel calculation, each $r$ is one complete panel
treated as exchangeable conditional on the common base Hamiltonian, not one
channel or one matrix entry.  Write $T^{(r)}$ for its physical
four-point trace and let $W^{(r,s)}$ denote the three-pairing Wick tensor whose
two second-moment factors are taken from distinct units $r\ne s$.  The
pair-split estimator is
\[
 \widehat{\mathcal K}
 =\frac{1}{R}\sum_{r=1}^{R}T^{(r)}
 -\frac{1}{R(R-1)}\sum_{r\ne s}W^{(r,s)}.
\]
Using distinct units in the quadratic products removes the leading
same-sample plug-in bias.  The ordered-pair form above is equivalent to an
unordered-pair average after symmetrizing $W^{(r,s)}$ and $W^{(s,r)}$.

For a registered scalar functional
$\widehat\theta=\mathcal L(\widehat{\mathcal K})$, delete unit $r$, recompute
the entire estimator as $\widehat\theta_{(-r)}$, and form the delete-one
pseudovalues
\[
 p_r=R\widehat\theta-(R-1)\widehat\theta_{(-r)},
 \qquad
 \overline p=R^{-1}\sum_r p_r.
\]
The reported jackknife standard error is
\[
 \operatorname{SE}_{\mathrm J}
 =\left[\frac{1}{R(R-1)}
 \sum_{r=1}^{R}(p_r-\overline p)^2\right]^{1/2}.
\]
This uncertainty is conditional on exchangeability of the declared units.  A
panel jackknife at one Hamiltonian measures variation within that tangent
construction; it cannot be relabeled as disorder or Hamiltonian uncertainty.

\subsection{Exact finite-population centering and frozen inference split}

The corrected complete-cumulant calculation uses 24 registered panels in
each case.  Centering is performed before the split by an exactly evaluated
mean of the finite tangent population, not by a sample average over the same
panels used for inference.  For a Moore--Read system with $N_\phi$ guiding
centers, the population consists of the $N_\phi^2$ torus anchors.  Linearity
of the protected response gives
\[
 \overline X_a=X\!\left(
 N_\phi^{-2}\sum_{x,y}G_{x,y}
 \right),
\]
and the uniform average of the traceless guiding-density generators vanishes.
For the local lattice-supercharge protocol, we enumerate the ordered
coordinate marginals generated by the registered randomized ordering and
average their projector responses.  This accounts for the nonzero population
mean of that discrete tangent construction.

After subtracting $\overline X_a$, panels $0,\ldots,11$ determine the Hermitian
channel Gram matrix.  Its positive inverse square root is frozen and applied
to the disjoint inference panels $12,\ldots,23$.  The training Gram after
whitening is required to agree with the $2\times2$ identity, and singular or
ill-conditioned channel support fails closed.  Interchanging the two halves
defines a reverse split, but its result is descriptive: it is not pooled and
does not carry a claim gate.

Let $\widehat\theta$ be the fixed scalar contraction of the centered complete
cumulant and $s_{\mathrm J}$ its delete-one jackknife error from the twelve
inference panels.  We form $t=\widehat\theta/s_{\mathrm J}$ and compare it with
a Student-$t_{11}$ reference.  Since five primary cases were registered, the
one-sided family level $0.05$ is implemented by the Bonferroni per-case level
$0.01$.  A case passes only if its one-sided $p$ value is below $0.01$ and the
corresponding one-sided lower bound is positive.  This gate is conditional on
exchangeability of complete panels at the fixed Hamiltonian and does not
estimate Hamiltonian-to-Hamiltonian fluctuations.

\subsection{Form-factor normalization and finite-rank references}

For a Hermitian matrix with eigenvalues $f_i$, our raw and connected form
factors are
\[
 K_{\mathrm{raw}}(\tau)=D^{-1}\mathbb E
 \left|\sum_{i=1}^{D}e^{-i\tau f_i}\right|^2,
 \qquad
 K_c(\tau)=K_{\mathrm{raw}}(\tau)
 -D^{-1}\left|\mathbb E\sum_{i=1}^{D}e^{-i\tau f_i}\right|^2.
\]
For exactly degenerate shifted energies, all $f_i=0$, so the disconnected
piece exhausts the raw form factor and $K_c=0$.  For curvature eigenvalues,
the same subtraction removes the one-point density and retains the connected
two-level correlation.

When the null is a finite complex-Jacobi point process on $[-1,1]$, its
continuous joint density has the form
\[
 p(\{f_i\})\propto
 \prod_{i<j}(f_i-f_j)^2
 \prod_i(1-f_i)^{\alpha}(1+f_i)^{\beta}.
\]
The finite-rank kernel is constructed from orthonormal Jacobi polynomials, and
the connected form factor is the Fourier transform of the corresponding
connected two-point function.  If the compression dimensions force exact
boundary atoms, those deterministic eigenvalues are separated before the
continuous kernel is compared.  This reference tests a specified finite-$D$
compression ensemble; agreement over a registered interval is not an
asymptotic universality statement.
